\begin{lemma}\label{lem:GrobnerQs}
If $w \in S_n$ has no obstruction of Type $1$, $2$, or $3$ and Conjecture \ref{conj:mainresult} holds for all permutations of smaller Coxeter length than that of $w$, then either $D_w = \Dom(w)$ or there is some lower outside corner $(i,j)$ of $D_w \setminus \Dom(w)$ corresponding to the variable $y$ in $Z_w$ so that $Q_{y,I_w}$ is CDG.
\end{lemma}
\begin{proof}
Fix a permutation $w \in S_n$ that has no obstruction of Type $1$, $2$, or $3$, and assume $D_w \neq \Dom(w)$.  First suppose that $D_w$ has some lower outside corner $(i,j)\notin \Dom(w)$ corresponding to the variable $y$ in $Z_w$ satisfying $\rank_w(i,j)+1 = \min\{m_1,m_2\}$, with notation as in Definition \ref{def:Qdef}.  

Write the CDG generators of $I$ as $\{yq_1+r_1, \ldots, yq_k+r_k, h_1, \ldots, h_\ell \}$, where $y$ does not divide any $q_i$, $r_i$, or $h_j$.  As discussed in Lemma \ref{lem:SchubertDel}, every $(\rank_w(i,j)+1)$-minor involving row $i$ and column $j$ has a term divisible by $y$, and so $\{q_1, \ldots, q_k\}$ generates the ideal of $\rank_w(i,j)$-minors in the submatrix of $Z_w$ strictly northwest of $y$.  Because the $Q_{y,I}$ is an ideal of maximal minors (after possible removing full rows or columns of $0$'s), the result follows from \cite[Theorem 4.2]{CDG15} or \cite[Proposition 5.4]{Boo11}.  

Alternatively, suppose that $D_w$ has no such lower outside corner, and fix any lower outside corner of $D_w$.  Because $Q_{y,I_w}$ depends only on $D_w$ weakly northwest of $(i,j)$, we may assume that $(i,j)$ is the only lower outside corner of $D_w$ and that $(1,j), (i,1) \notin \Dom(w)$.  With this assumption, $Q_{y,I_w} +(z_{a,b} \mid (a,b) \in \Dom(w))= C_{y,I_w}$.  Because the $z_{a,b}$ with $(a,b) \in \Dom(w)$ are indeterminates that do not divide any term of any CDG generator of $Q_{y,I_w}$, it follows that $Q_{y,I_w}$ is CDG if and only if $C_{y,I_w}$ is CDG.  

If $\Dom(w) = \emptyset$, then, because $w$ has no obstruction of Type $1$, $w$ must be vexillary and so $I_w$ is CDG by \cite[Theorem 3.8]{KMY09}.   It then follows from  \cite[Theorem 2.1(a)]{KMY09} that $C_{y,I_w}$ is CDG as well.  

Now suppose $\Dom(w) \neq \emptyset$.  The assumptions that $(1,j), (i,1) \notin \Dom(w)$ imply that there is at least one element of $\Dom(w) \cap \Ess(w)$ northwest of $(i,j)$.  Choose the southernmost such element and label it $(r,s)$ and the easternmost element and label it $(r',s')$.  Because $w$ has no obstruction of Type $1$, all essential boxes of $D_w \setminus \Dom(w)$ must be in either row $i$ or column $j$.  Because $\rank_w(i,j)+1 < \min\{i,j\}$, there must be at least one essential box in row $i$ and at least one in column $j$ aside from $(i,j)$. Then because $w$ has no obstruction of Type $2$, there are no elements of $\Ess(w)$ north of row $i$ and south of row $r$ or west of column $j$ and east of column $s'$.  

We will argue directly in this case that for each pair $q_a$ and $h_b$, their $s$-polynomial has a Gr\"obner reduction by the CDG generators of $Q_{y,I_w}$.  (For an example illustrating this process, see Example \ref{ex:Spoly} below.)  Choose such a $q_a$ and $h_b$.  Because the CDG generators of $N_{y,I_w}$ form a Gr\"obner basis by induction on the Coxeter length of $w$ and Corollary \ref{cor:GrobnerDel}, we may assume that $q_a \notin (h_1, \ldots, h_\ell)$.  

Because $Q_{y,I_w}$ involves only indeterminates of $Z_w$ northwest of $z_{i,j}$, we will work within the submatrix of $Z_w$ northwest of $z_{i,j}$, which we will call $Y_w$.  Suppose that $h_b$ is determined by the essential box at $(i',j)$ for some $i'<i$.  (The case of $(i,j')$ with $j'<j$ will follow by symmetry.)  We consider two cases.  

Case 1: Suppose that $(i',j-1) \in D_w$.  Then, because $w$ has no obstruction of Type $1$, there can be no element of $\Dom(w) \cap \Ess(w)$ northwest of $(i',j)$.  In particular, $r' \geq i'$, and $h_b$ is a $(\rank_w(i',j)+1)$-minor in the submatrix of $Y_w$ formed of its final $j-s'$ columns, which is a generic matrix and which we will call $Y'_w$.  Suppose that there are $t$ columns determining $q_a$ strictly east of column $s'$ and $\rank_w(i,j)-t$ columns determining $q_a$ weakly west of column $s'$.  Express $q_a$ as a sum of products of $t$-minors and $(\rank_w(i,j)-t)$-minors corresponding to this subdivision.  For any fixed set of rows of $Y'_w$, the set of $(\rank_w(i',j)+1)$-minors weakly northwest of $(i',j)$ together with the $t$-minors strictly northwest of $(i,j)$ in the submatrix of $Y'_w$ including only those specified rows forms a Gr\"obner basis for the ideal they generate because it is a mixed ladder determinantal ideal \cite[Theorem 1.10]{Gor07}.  

Choose the $t$-minor $\varepsilon_1$ from the eastern $t$ columns determining $q_a$ and the $(\rank_w(i,j)-t)$-minor $\varepsilon_2$ from the remaining columns satisfying $LT(\varepsilon_1) \cdot LT(\varepsilon_2) = LT(q_a)$.  Because $\varepsilon_1$ and $h_b$ belong to the ideal of $t$-minors weakly north of their southernmost entries together with the $(\rank_w(i',j)+1)$-minors northwest of $(i',j)$ in $Y'_w$, their $s$-polynomial $s(\varepsilon_1,h_b)$ has a Gr\"obner reduction $s(\varepsilon_1,h_b) = \sum \alpha_c \delta_c $ by the natural generators of that ideal.  For each $\delta_c \in N_{y,I_w}$, set $\widehat{\delta_c} = LT(\varepsilon_2) \cdot \delta_c$, and for each $\delta_c$ involving some row south of $i'$, set $\widehat{\delta_c}$ to be (up to sign) the determinant of the augmentation of the matrix determining $\delta_c$ by the rows and columns determining $\varepsilon_2$ (with sign chosen so that $LT(\delta_c)$ and $LT(\widehat{\delta_c})$ share a sign).  Let $s(q_a,h_b)$ denote the $s$-polynomial of $q_a$ and $h_b$.  

We claim that $s(q_a,h_b)-\sum \alpha_c \widehat{\delta_c} \in N_{y,I}$.  It is clear that $s(q_a,h_b)-\sum \alpha_c \widehat{\delta_c}$ contains a $LT(\varepsilon_2)$-multiple of $s(\varepsilon_1,h_b)-\sum \alpha_c \delta_c$, which is $0$ because $s(\varepsilon_1,h_b)-\sum \alpha_c \delta_c$ is.  Fix any non-leading term $\mu$ of $\varepsilon_2$, and write $s(q_a,h_b)-\sum \alpha_c \widehat{\delta_c} = \mu \tilde{s}+\tilde{\tilde{s}}$ where $\mu$ does not divide any term of $\tilde{\tilde{s}}$.  For each column involved in $\varepsilon_2$, whenever a different variable from that column divides $\mu$ and $LT(\varepsilon_2)$, replace in $\sum \alpha_c \widehat{\delta_c}$ the variables in the row of the divisor of $LT(\varepsilon_2)$ with the variables in the same columns from the row of the corresponding divisor of $\mu$, which we note gives an expression of $\tilde{s}$ in terms of the natural generators of the ideal of $(\rank_w(i',j)+1)$-minors northwest of $(i',j)$, each of which is a CDG generator of $N_{y,I}$.

Now because $LT(\alpha_c\delta_c) \leq LT(s(\varepsilon_1,h_b))$ and $LT(\alpha_c \widehat{\delta_c}) = LT(\varepsilon_2) \cdot LT(\alpha_c\delta_c)$ for each $c$ and because $LT(s(q_a,h_b)) = LT(\varepsilon_2) \cdot LT(s(\varepsilon_1,h_b))$, subtracting each $\alpha_c \widehat{\delta_c}$ is a valid step in a Gr\"obner reduction of $s(q_a,h_b)$ by the generators of $Q_{y,I_w}$.  The fact that the CDG generators of $N_{y,I_w}$, each of which is a CDG generator of $Q_{y,I_w}$, form a Gr\"obner basis for the ideal they generate implies that $s(q_a,h_b)-\sum \alpha_c \widehat{\delta_c} \in N_{y,I_w}$ has a reduction in terms of those generators and so that $s(q_a,h_b)$ has a reduction by the CDG generators of $Q_{y,I_w}$.

Case 2: Suppose $(i',j-1) \notin D_w$.  Then there can be no $j''<j$ with $(i',j'') \in D_w\setminus \Dom(w)$ because $w$ has no obstruction of Type $3$.  Hence, $\rank_w(i',j)+1 = \min\{j-k \mid (i',k) \in \Dom(w)\}$, and so the $(\rank_w(i',j)+1)$-minors northwest of $(i',j)$ are the maximal minors of the submatrix of $Z_w$ northwest of $(i',j)$ after removing complete rows or columns of $0$'s.  Then the argument is similar to the first case but uses, instead of results on ladder determinantal ideals in a generic matrix, the fact that the maximal minors of matrices of indeterminates and $0$'s form a Gr\"obner basis by \cite[Theorem 4.2]{CDG15} or \cite[Proposition 5.4]{Boo11}.  
\end{proof}
